\documentclass{article}
\usepackage{graphicx,amsfonts,amsmath,titlesec,geometry} % Required for inserting images
\renewcommand{\thesection}{(\arabic{section})}
\renewcommand{\thesubsection}{\alph{subsection})}
\renewcommand{\thesubsubsection}{\arabic{subsubsection}.}
% \geometry{margin=20mm}
\title{Introduction to Coding Theory HW 3}
\author{Alex}

\newcommand{\wt}{\text{wt }}


\begin{document}
\maketitle
\section{}
    \subsection{}
        We have \(C\) is given by \(F(X,Y,Z) = X^{q+1}-Y^qZ-YZ^q=0\).

        First, consider case \(Z=0, Y=1\).
        We have \(F_X=(q+1)X^q=0\implies X^q = 0\).
        We have \(F_Z=-Y^q-qYZ^{q-1}=-Y^q=-(1)^q\ne 0\) so we never have
        a singular point in the affine plane.

        Consider next case \(Z=1\). We then have \(F_Z=-Y^q-qYZ^{q-1}=
        -Y^q-qY\) and similarly \(F_Y=-qY^{q-1}-Z^q\). At any singular point
        We have already \(X^q=0\) so we have \(F(X,Y,1)=X^qX-Y^q-Y=0\rightarrow
        -Y^q-Y=0\). So, we have \(Y(Y^q+1)=0\). If we assume \(Y=0\) then we 
        ger $F_Y=0-Z^q=1$, which can not occur, so we have \(Y^q=-1\). Though
        we then have \(F_Z=1-qY=0\) so we have \(Y=\frac1q\) so we have 
        \(F\left(X,\frac1q,1\right)= 0+1-\frac1q=0\rightarrow q = 1\), though
        we have $q$ is prime so this is a contradiction, so there are no
        singular points on $C$ so $C$ is nonsingular.
    \subsection{}
        First consider affine points with \(Z=1\). These then must satisfy
        \(X^3=Y^2-Y\). We have $X^3\in\{0,1,2^3=0,3^3=3\}$ and similarly
        $Y^2-Y\in\{0,0,2^2-2=-$
\end{document}